\documentclass[10pt,letterpaper]{article}
\usepackage[margin=0.78in,headheight=17pt]{geometry}
\usepackage[T1]{fontenc}
\usepackage[utf8]{inputenc}
\usepackage{lmodern,microtype,parskip}
\usepackage{amsmath,amssymb}
\usepackage{xcolor,booktabs,longtable,tabularx,array,enumitem}
\usepackage[most]{tcolorbox}
\usepackage{fancyhdr,lastpage,titlesec,hyperref,listings}
\usepackage{tikz}
\usetikzlibrary{arrows.meta,positioning,shapes.geometric}
\definecolor{navy}{HTML}{17324D}
\definecolor{bluegray}{HTML}{43657D}
\definecolor{lightblue}{HTML}{EAF2F8}
\definecolor{linegray}{HTML}{CBD5DF}
\definecolor{softgray}{HTML}{F3F5F7}
\definecolor{accent}{HTML}{8E3B46}
\hypersetup{colorlinks=true,linkcolor=navy,urlcolor=bluegray,pdfborder={0 0 0}}
\urlstyle{same}
\titleformat{\section}{\Large\bfseries\color{navy}}{\thesection}{0.6em}{}
\titleformat{\subsection}{\large\bfseries\color{bluegray}}{\thesubsection}{0.5em}{}
\setlist[itemize]{leftmargin=*,itemsep=3pt,topsep=3pt}
\setlist[enumerate]{leftmargin=*,itemsep=3pt,topsep=3pt}
\setlength{\emergencystretch}{3em}
\renewcommand{\arraystretch}{1.22}
\setlength{\tabcolsep}{4pt}
\pagestyle{fancy}\fancyhf{}
\lhead{Secure Programming \& SAST Engineering}
\rhead{Engineering Documentation Suite}
\cfoot{\thepage\ / \pageref{LastPage}}
\newcolumntype{Y}{>{\raggedright\arraybackslash}X}
\newcommand{\checkbox}{\(\square\)}
\newcommand{\blankline}{\noindent\rule{\linewidth}{0.35pt}\par\vspace{4pt}}
\newcommand{\writearea}[1]{\par\begin{tcolorbox}[colback=white,colframe=linegray,boxrule=0.5pt,height=#1,sharp corners]\mbox{}\end{tcolorbox}\par}
\newtcolorbox{focusbox}[1]{colback=lightblue,colframe=bluegray,title=\textbf{#1},fonttitle=\bfseries,breakable}
\newtcolorbox{warnbox}[1]{colback=softgray,colframe=accent,title=\textbf{#1},fonttitle=\bfseries,breakable}
\lstset{basicstyle=\ttfamily\footnotesize,breaklines=true,frame=single,rulecolor=\color{linegray},backgroundcolor=\color{softgray},showstringspaces=false,columns=fullflexible,keepspaces=true,upquote=true}

\begin{document}
\begin{titlepage}
\vspace*{1.0in}
{\color{bluegray}\Large Analysis engineering lab}\\[1.0em]
{\color{navy}\Huge\bfseries CFG, Dataflow and Taint Exercise}\\[0.7em]
{\Large A finite-lattice fixed-point analysis with reviewable traces}\\[2.0em]
\rule{\linewidth}{1.0pt}
\vfill
\textbf{Source discipline:} Book-derived foundations are identified separately from current engineering extensions.\\
\textbf{Edition:} 29 October 2026
\end{titlepage}
\tableofcontents
\clearpage
\section{Goal}
Build a small forward dataflow analysis on a directed control-flow graph (CFG), derive fixed-point states and document known limitations. This is an analysis exercise, not a production-grade security scanner.
\section{Input program}
\begin{lstlisting}
x = request.input      # node 1: source
if admin_mode:          # node 2
    x = allowlist(x)    # node 3
else:
    x = x               # node 4
execute_sql(x)         # node 5: sink
\end{lstlisting}
\section{Model and iterative algorithm}
Represent each node's state as a set of possibly tainted variables. Join predecessors with union. Transfer: assignment from a tainted value taints the target; sanitization clears only when its contract is correct for the destination sink. To teach conservative behavior, model \texttt{allowlist} as a known sanitizer for a very narrow input grammar; treat all unknown transformations as propagating taint.
\begin{lstlisting}
worklist = all_nodes
IN[n] = empty_set; OUT[n] = empty_set
while worklist:
    n = worklist.pop()
    new_in = union(OUT[p] for p in predecessors[n])
    new_out = transfer(n, new_in)
    if (new_in, new_out) != (IN[n], OUT[n]):
        IN[n], OUT[n] = new_in, new_out
        worklist.add_all(successors[n])
\end{lstlisting}
\section{Tasks}
\begin{enumerate}
\item Draw explicit CFG edges, including both branches and the join at node 5.
\item Fill iterations 0, 1, 2, ... until no state changes; record IN/OUT at each node.
\item Explain why sanitizing only one branch does not prove the sink safe.
\item Introduce an alias (\texttt{y = x}) and a helper function call; explain missing alias or call modeling.
\item Define two false-positive fixtures and one false-negative fixture for your simplified analyzer.
\end{enumerate}
\section{Exit artifacts}
Graph, transfer table, convergence log, sink trace, limitations note and corrected control invariant. Passing requires detecting the unclean else-path and naming a precision limitation.
\end{document}
